\documentclass[10pt,a4paper]{amsart}
\usepackage{amsmath,amssymb,amsthm,graphicx}
% Retained minimal preamble for the Stacks Project chapter more-algebra.tex (More on Algebra).
% The chapter file pinned-original-source.tex carries only \input{preamble} on its first line, so no
% preamble text lives inside the chapter itself.  The Stacks Project uses one single shared preamble for
% every chapter; the snapshot stages it identically as sources/stacks-chow/preamble.tex and
% sources/stacks-etale-cohomology/preamble.tex.  The block below is copied mechanically, line for line,
% from sources/stacks-chow/preamble.tex lines 14-18 (the Xy-pic setup) and 154-217 (the native theorem
% declarations together with the mathematical macros).  Nothing was added, removed or redefined; the
% wrapper class and the amsmath/amssymb/amsthm/graphicx base packages are supplied by the assembler.
% Chapter: more-algebra.tex -- More on Algebra.
% For commutative diagrams we use Xy-pic
\usepackage[all]{xy}

% We use 2cell for 2-commutative diagrams.
\xyoption{2cell}

% Theorem environments.
%
\theoremstyle{plain}
\newtheorem{theorem}[subsection]{Theorem}
\newtheorem{proposition}[subsection]{Proposition}
\newtheorem{lemma}[subsection]{Lemma}

\theoremstyle{definition}
\newtheorem{definition}[subsection]{Definition}
\newtheorem{example}[subsection]{Example}
\newtheorem{exercise}[subsection]{Exercise}
\newtheorem{situation}[subsection]{Situation}

\theoremstyle{remark}
\newtheorem{remark}[subsection]{Remark}
\newtheorem{remarks}[subsection]{Remarks}

\numberwithin{equation}{subsection}

% Macros
%
\def\lim{\mathop{\mathrm{lim}}\nolimits}
\def\colim{\mathop{\mathrm{colim}}\nolimits}
\def\Spec{\mathop{\mathrm{Spec}}}
\def\Hom{\mathop{\mathrm{Hom}}\nolimits}
\def\Ext{\mathop{\mathrm{Ext}}\nolimits}
\def\SheafHom{\mathop{\mathcal{H}\!\mathit{om}}\nolimits}
\def\SheafExt{\mathop{\mathcal{E}\!\mathit{xt}}\nolimits}
\def\Sch{\mathit{Sch}}
\def\Mor{\mathop{\mathrm{Mor}}\nolimits}
\def\Ob{\mathop{\mathrm{Ob}}\nolimits}
\def\Sh{\mathop{\mathit{Sh}}\nolimits}
\def\NL{\mathop{N\!L}\nolimits}
\def\CH{\mathop{\mathrm{CH}}\nolimits}
\def\proetale{{pro\text{-}\acute{e}tale}}
\def\etale{{\acute{e}tale}}
\def\QCoh{\mathit{QCoh}}
\def\Ker{\mathop{\mathrm{Ker}}}
\def\Im{\mathop{\mathrm{Im}}}
\def\Coker{\mathop{\mathrm{Coker}}}
\def\Coim{\mathop{\mathrm{Coim}}}

% Boxtimes
%
\DeclareMathSymbol{\boxtimes}{\mathbin}{AMSa}{"02}

%
% Macros for moduli stacks/spaces
%
\def\QCohstack{\mathcal{QC}\!\mathit{oh}}
\def\Cohstack{\mathcal{C}\!\mathit{oh}}
\def\Spacesstack{\mathcal{S}\!\mathit{paces}}
\def\Quotfunctor{\mathrm{Quot}}
\def\Hilbfunctor{\mathrm{Hilb}}
\def\Curvesstack{\mathcal{C}\!\mathit{urves}}
\def\Polarizedstack{\mathcal{P}\!\mathit{olarized}}
\def\Complexesstack{\mathcal{C}\!\mathit{omplexes}}
% \Pic is the operator that assigns to X its picard group, usage \Pic(X)
% \Picardstack_{X/B} denotes the Picard stack of X over B
% \Picardfunctor_{X/B} denotes the Picard functor of X over B
\def\Pic{\mathop{\mathrm{Pic}}\nolimits}
\def\Picardstack{\mathcal{P}\!\mathit{ic}}
\def\Picardfunctor{\mathrm{Pic}}
\def\Deformationcategory{\mathcal{D}\!\mathit{ef}}

% --- appended 2026-10-12 (Chief): the project own definitions of the reference/slogan/history
% display environments, copied mechanically from the shared Stacks preamble lines 8-12
% (\usepackage{verbatim} plus the three \newenvironment lines, which the chapter preamble keeps
% verbatim from the project preamble). Nothing added, removed or redefined: these environments
% carry no mathematical content and the project itself typesets them as comments. Without these
% three lines a counted statement whose own text uses them fails with Environment slogan undefined.
\usepackage{verbatim}
\newenvironment{reference}{\comment}{\endcomment}
\newenvironment{slogan}{\comment}{\endcomment}
\newenvironment{history}{\comment}{\endcomment}
\usepackage{cleveref}
\crefname{originalremark}{original remark}{original remarks}
\Crefname{originalremark}{Original remark}{Original remarks}
\makeatletter
\newlabel{remark-construction}{{09EM}{source}{}{sourcepointer-1}{}}
\newlabel{remark-construction@cref}{{[originalremark][0][]09EM}{[1][0][]source}{}{}{}}
\makeatother
\crefname{originallemma}{original lemma}{original lemmas}
\Crefname{originallemma}{Original lemma}{Original lemmas}
\makeatletter
\newlabel{lemma-extension-dvrs-formally-smooth}{{09E7}{source}{}{sourcepointer-2}{}}
\newlabel{lemma-extension-dvrs-formally-smooth@cref}{{[originallemma][0][]09E7}{[1][0][]source}{}{}{}}
\makeatother
\title{(Weak) solutions descend along the intermediate extension of discrete valuation rings and compose}
\author{The Stacks Project Authors}
\date{stacks-project:more-algebra.tex:local-snapshot-20261012; GNU FDL 1.2 or later}
\begin{document}
\maketitle
\setcounter{subsection}{1}
\noindent\textit{Editorial note.} All delivered mathematical text is the author's own, copied from the pinned Stacks Project chapter file more-algebra.tex; no statement or proof step was rewritten, reordered or removed. The counted statement is the labeled lemma environment at source lines 34516-34527 (canonical tag 09ES) and the counted proof is the proof environment at source lines 34529-34552, which immediately follows it and is its own proof.\par\medskip
\noindent\textit{Editorial note.} The proof ends with the sentence that a similar argument works for the second assertion; that sentence is the author's own and is preserved verbatim inside the counted span, the second case being the same localization argument applied to the compositional hypothesis.\par\medskip
\noindent\textit{Editorial note.} The two references are to earlier results of this chapter, the construction remark and the formally smooth criterion; both are navigation stubs with their roles recorded in the reasons.\par\medskip
\noindent\textit{Editorial note.} The preamble span runs to line 90, the current end of the retained preamble, which includes the appended environment definitions; the counted text does not use them, and the span is taken whole so that the file version is unambiguous.\par\medskip
\noindent\textit{Editorial note.} The source URL is the source-level Stacks Project URL, as required when more than one card is produced from one source; the per-claim tag is carried by claim\_key and by this note. The counted statement and proof cite no bibliography entry, so this card needs no bibliography.\par\medskip
\noindent\textit{Allowed standard prerequisites (editorial).}\begin{itemize}
\item The reduced base change of a finite compositum is a finite product of finite extensions and the corresponding map of integral closures is a product of the maps considered in the construction remark, hence surjective on spectra. invoked at proof lines 34529-34538. Provides a maximal ideal of the larger integral closure lying over a chosen prime of the smaller one, which is the point at which the comparison is made.
\item Weak unramifiedness for a composite of extensions of discrete valuation rings implies weak unramifiedness of the intermediate extension, and the same holds for separability of the residue field extensions. used at proof lines 34538-34547. Yields assertion (1) by comparing the composite with its first factor.
\item The equivalence between weak unramifiedness with separable residue extension and formal smoothness in the adic topology, which is the definition of a weak solution. invoked at proof lines 34547-34549. Translates the conclusion about weak unramifiedness and separable residue fields back into the statement about weak solutions.
\end{itemize}\medskip
\ifdefined\hypertarget\hypertarget{sourcepointer-1}{}\fi
\noindent\textit{Source reference.} Original remark 09EM: sources/stacks-more-algebra/pinned-original-source.tex, source lines 33882--33882. The reference identifies the map between the integral closures as a finite product of the standard maps, used at proof lines 34533-34538 to get surjectivity on spectra; the localization argument built on it is inside the counted span.\par\medskip
\ifdefined\hypertarget\hypertarget{sourcepointer-2}{}\fi
\noindent\textit{Source reference.} Original lemma 09E7: sources/stacks-more-algebra/pinned-original-source.tex, source lines 32989--32989. The reference is invoked at proof lines 34547-34549 to translate the conclusions about weak unramifiedness and separable residue extensions into the statement about (weak) solutions; it is the last step and everything leading to it is inside the counted span.\par\medskip
\begin{lemma}
\label{lemma-solutions-go-down}
Let $A \to B \to C$ be extensions of discrete valuation rings with fraction
fields $K \subset L \subset M$. Let $K_1/K$ be a finite extension.
\begin{enumerate}
\item If $K_1$ is a (weak) solution for $A \to C$, then $K_1$ is a (weak)
solution for $A \to B$.
\item If $K_1$ is a (weak) solution for $A \to B$ and
$L_1 = (L \otimes_K K_1)_{red}$ is a product of fields which are
(weak) solutions for $B \to C$, then $K_1$ is a (weak) solution for $A \to C$.
\end{enumerate}
\end{lemma}

\begin{proof}
Let $L_1 = (L \otimes_K K_1)_{red}$ and $M_1 = (M \otimes_K K_1)_{red}$
and let $B_1 \subset L_1$ and $C_1 \subset M_1$ be the integral closure
of $B$ and $C$. Note that $M_1 = (M \otimes_L L_1)_{red}$ and that $L_1$
is a (nonempty) finite product of finite extensions of $L$. Hence the
ring map $B_1 \to C_1$ is a finite product of ring maps of the form discussed
in Remark \ref{remark-construction}. In particular, the map
$\Spec(C_1) \to \Spec(B_1)$ is surjective. Choose a maximal ideal
$\mathfrak m \subset C_1$ and consider the extensions of discrete
valuation rings
$$
(A_1)_{A_1 \cap \mathfrak m} \to
(B_1)_{B_1 \cap \mathfrak m} \to
(C_1)_\mathfrak m
$$
If the composition is weakly unramified, so is the map
$(A_1)_{A_1 \cap \mathfrak m} \to (B_1)_{B_1 \cap \mathfrak m}$.
If the residue field extension
$\kappa_{A_1 \cap \mathfrak m} \to \kappa_\mathfrak m$ is separable,
so is the subextension
$\kappa_{A_1 \cap \mathfrak m} \to \kappa_{B_1 \cap \mathfrak m}$.
Taking into account Lemma \ref{lemma-extension-dvrs-formally-smooth}
this proves (1). A similar argument works for (2).
\end{proof}

\end{document}
